% ============================================================
% SEBC Dushigikirane — Manuel d'Utilisation
% Guide complet pour les utilisateurs
% ============================================================
\documentclass[12pt,a4paper]{article}
\usepackage[utf8]{inputenc}
\usepackage[T1]{fontenc}
\usepackage[french]{babel}
\usepackage{geometry}
\geometry{margin=2.5cm}
\usepackage{graphicx}
\usepackage{xcolor}
\usepackage{hyperref}
\usepackage{fancyhdr}
\usepackage{titlesec}
\usepackage{enumitem}
\usepackage{tabularx}
\usepackage{booktabs}
\usepackage{tcolorbox}
\usepackage{fontawesome5}
\usepackage{mdframed}

% ═══ Couleurs SEBC ═══
\definecolor{sebcblue}{HTML}{1A3A5C}
\definecolor{sebcblulight}{HTML}{2C5F8A}
\definecolor{sebcred}{HTML}{C8102E}
\definecolor{sebcteal}{HTML}{0D9488}
\definecolor{sebcgray}{HTML}{64748B}
\definecolor{lightbg}{HTML}{F0F4FF}
\definecolor{warnbg}{HTML}{FEF3C7}
\definecolor{successbg}{HTML}{D1FAE5}
\definecolor{infobg}{HTML}{DBEAFE}

% ═══ Styles ═══
\hypersetup{
    colorlinks=true,
    linkcolor=sebcblue,
    urlcolor=sebcblulight,
}

\titleformat{\section}{\Large\bfseries\color{sebcblue}}{}{0em}{}[\titlerule]
\titleformat{\subsection}{\large\bfseries\color{sebcblulight}}{}{0em}{}
\titleformat{\subsubsection}{\normalsize\bfseries\color{sebcgray}}{}{0em}{}

\pagestyle{fancy}
\fancyhf{}
\fancyhead[L]{\small\color{sebcgray}SEBC — Manuel d'Utilisation}
\fancyhead[R]{\small\color{sebcgray}\thepage}
\fancyfoot[C]{\small\color{sebcgray}Guide utilisateur — \today}
\renewcommand{\headrulewidth}{0.5pt}
\renewcommand{\footrulewidth}{0.3pt}

\tcbuselibrary{skins,breakable}

% Boîte étape
\newtcolorbox{etapebox}[1][]{
    colback=lightbg,
    colframe=sebcblue,
    fonttitle=\bfseries,
    title=#1,
    rounded corners,
    boxrule=0.5pt,
    left=8pt, right=8pt, top=6pt, bottom=6pt,
}

% Boîte astuce
\newtcolorbox{astucebox}[1][]{
    colback=successbg,
    colframe=sebcteal,
    fonttitle=\bfseries\color{sebcteal},
    title={\faLightbulb~#1},
    rounded corners,
    boxrule=0.5pt,
    left=8pt, right=8pt,
}

% Boîte attention
\newtcolorbox{attentionbox}[1][]{
    colback=warnbg,
    colframe=orange,
    fonttitle=\bfseries\color{orange!80!black},
    title={\faExclamationTriangle~#1},
    rounded corners,
    boxrule=0.5pt,
    left=8pt, right=8pt,
}

% Boîte info
\newtcolorbox{infobox}[1][]{
    colback=infobg,
    colframe=sebcblue,
    fonttitle=\bfseries\color{sebcblue},
    title={\faInfoCircle~#1},
    rounded corners,
    boxrule=0.5pt,
    left=8pt, right=8pt,
}

\begin{document}

% ═══════════════════════════════════════════════════════════
% PAGE DE TITRE
% ═══════════════════════════════════════════════════════════
\begin{titlepage}
\begin{center}
\vspace*{2cm}

{\Huge\bfseries\color{sebcblue} SEBC Dushigikirane}\\[0.5cm]
{\Large\color{sebcblulight} Soutien Entre Burundais du Canada}\\[2cm]

\rule{\textwidth}{2pt}\\[0.8cm]
{\LARGE\bfseries Manuel d'Utilisation}\\[0.3cm]
{\large\color{sebcgray} Guide complet pour les membres}\\[0.5cm]
\rule{\textwidth}{2pt}\\[2cm]

{\large\color{sebcgray}
\faGlobe\ \url{https://sebc-dushigikirane.pro}\\[1cm]
}

\begin{tabular}{ll}
\textbf{Version} & 1.0 \\
\textbf{Date} & \today \\
\textbf{Public cible} & Tous les membres SEBC \\
\end{tabular}

\vfill
{\small\color{sebcgray} SEBC Dushigikirane — Entraidons-nous}
\end{center}
\end{titlepage}

% ═══════════════════════════════════════════════════════════
\tableofcontents
\newpage

% ═══════════════════════════════════════════════════════════
\section{Accéder à la Plateforme}
% ═══════════════════════════════════════════════════════════

\subsection{Adresse du site}

La plateforme SEBC est accessible à l'adresse suivante :

\begin{center}
\large\textbf{\url{https://sebc-dushigikirane.pro}}
\end{center}

\begin{infobox}[Navigateurs compatibles]
La plateforme fonctionne sur tous les navigateurs modernes :
\begin{itemize}[leftmargin=1.5em]
    \item Google Chrome (recommandé)
    \item Mozilla Firefox
    \item Microsoft Edge
    \item Safari (Mac/iPhone)
\end{itemize}
La plateforme est également optimisée pour les \textbf{téléphones mobiles} et les \textbf{tablettes}.
\end{infobox}

% ═══════════════════════════════════════════════════════════
\section{Devenir Membre (Candidature)}
% ═══════════════════════════════════════════════════════════

Pour rejoindre SEBC, vous devez passer par un processus d'adhésion en deux étapes. Un \textbf{parrain} (membre existant) est obligatoire.

\subsection{Étape 1 — Vérification du Parrain}

\begin{etapebox}[Procédure]
\begin{enumerate}[leftmargin=1.5em]
    \item Rendez-vous sur \textbf{\url{https://sebc-dushigikirane.pro/candidature/}}
    \item Saisissez l'\textbf{adresse email} ou le \textbf{numéro de téléphone} de votre parrain
    \item Cliquez sur \textbf{Vérifier}
    \item Si le parrain est trouvé, son email s'affiche (par confidentialité, aucune autre information n'est partagée)
    \item Cliquez sur \textbf{Continuer} pour passer au formulaire d'identification
\end{enumerate}
\end{etapebox}

\begin{attentionbox}[Important]
Votre parrain doit être un membre \textbf{déjà enregistré et approuvé} dans la base de données SEBC. Si vous n'avez pas de parrain, contactez un membre existant pour qu'il vous parraine.
\end{attentionbox}

\subsection{Étape 2 — Identification}

\begin{etapebox}[Formulaire d'identification]
Remplissez les champs suivants :
\begin{itemize}[leftmargin=1.5em]
    \item \textbf{Nom} et \textbf{Prénom}
    \item \textbf{Email} (unique — sera votre identifiant de connexion)
    \item \textbf{Téléphone WhatsApp} (obligatoire)
    \item \textbf{Téléphone Canada} (optionnel)
    \item \textbf{Province d'origine} (Burundi)
    \item \textbf{Ville de résidence} et \textbf{Pays de résidence}
    \item \textbf{Noms des parents}, conjoint, enfants (optionnels)
    \item \textbf{Nombre de personnes en charge}
\end{itemize}
\end{etapebox}

\subsection{Étape 3 — Vérification Email}

\begin{enumerate}
    \item Après soumission du formulaire, un \textbf{code OTP à 6 chiffres} est envoyé à votre email
    \item Saisissez le code dans le champ prévu
    \item Le code est valide pendant \textbf{15 minutes}
\end{enumerate}

\begin{astucebox}[Astuce]
Vérifiez votre dossier \textbf{Spam/Courrier indésirable} si vous ne recevez pas le code. L'expéditeur est \texttt{link2ictgroup@gmail.com}.
\end{astucebox}

\subsection{Étape 4 — Création du mot de passe}

\begin{enumerate}
    \item Après la vérification, vous êtes invité à \textbf{créer un mot de passe}
    \item Le mot de passe doit être confirmé (deux saisies identiques)
    \item Votre compte est maintenant créé avec le statut \textbf{En attente}
\end{enumerate}

\subsection{Étape 5 — Validation du parrain}

\begin{enumerate}
    \item Votre parrain reçoit un \textbf{email de notification}
    \item Il doit se connecter à la plateforme pour \textbf{valider votre parrainage}
    \item Vous pouvez \textbf{relancer} votre parrain depuis votre espace si nécessaire
\end{enumerate}

% ═══════════════════════════════════════════════════════════
\section{Se Connecter}
% ═══════════════════════════════════════════════════════════

\begin{etapebox}[Procédure de connexion]
\begin{enumerate}[leftmargin=1.5em]
    \item Rendez-vous sur \textbf{\url{https://sebc-dushigikirane.pro/login/}}
    \item Saisissez votre \textbf{adresse email}
    \item Cliquez sur \textbf{Continuer}
    \item Si votre email est reconnu, saisissez votre \textbf{mot de passe}
    \item Cliquez sur \textbf{Se connecter}
\end{enumerate}
\end{etapebox}

Après connexion réussie, vous êtes redirigé vers votre \textbf{Espace Membre}.

% ═══════════════════════════════════════════════════════════
\section{Navigation — La Sidebar}
% ═══════════════════════════════════════════════════════════

La \textbf{sidebar} (barre latérale gauche) affiche les modules auxquels vous avez accès :

\begin{tabularx}{\textwidth}{clX}
\toprule
\textbf{Icône} & \textbf{Module} & \textbf{Description} \\
\midrule
\faUser & Mon Espace & Votre tableau de bord personnel \\
\faBullhorn & Communication & Messagerie et vidéoconférence \\
\faCogs & Administration & Gestion de l'association (administrateurs uniquement) \\
\bottomrule
\end{tabularx}

\begin{infobox}[Modules selon votre rôle]
Les modules visibles dans la sidebar dépendent de votre \textbf{rôle} dans l'association. Un membre simple verra « Mon Espace » et « Communication », tandis qu'un administrateur verra tous les modules.
\end{infobox}

Sur \textbf{mobile}, la sidebar est masquée par défaut. Cliquez sur le bouton menu (\faBar) dans la barre supérieure pour l'afficher.

% ═══════════════════════════════════════════════════════════
\section{Mon Espace (Module Membres)}
% ═══════════════════════════════════════════════════════════

\subsection{Tableau de bord}

Votre espace personnel affiche :
\begin{itemize}
    \item Vos \textbf{informations personnelles} (modifiables)
    \item Le \textbf{statut de votre compte} (en attente, approuvé, etc.)
    \item Votre \textbf{cellule} d'appartenance
    \item La liste de vos \textbf{ayants droits} déclarés
\end{itemize}

\subsection{Gestion des Ayants Droits}

Les ayants droits sont les personnes que vous déclarez comme bénéficiaires en cas de soutien.

\begin{etapebox}[Ajouter un ayant droit]
\begin{enumerate}[leftmargin=1.5em]
    \item Cliquez sur le bouton \textbf{+ Ajouter un ayant droit}
    \item Remplissez : nom, prénom, type de lien (père, mère, conjoint, enfant, etc.)
    \item Joignez un \textbf{document justificatif} si requis
    \item Cliquez sur \textbf{Enregistrer}
    \item L'ayant droit apparaît avec le statut \textbf{En attente d'approbation}
\end{enumerate}
\end{etapebox}

\subsection{Modification du profil}

Vous pouvez mettre à jour vos informations personnelles :
\begin{itemize}
    \item Adresse, ville, pays de résidence
    \item Numéros de téléphone
    \item Informations familiales
\end{itemize}

\begin{attentionbox}[Attention]
L'\textbf{adresse email} ne peut pas être modifiée car elle sert d'identifiant unique. Contactez un administrateur si nécessaire.
\end{attentionbox}

\subsection{Parrainage}

Si vous êtes parrain d'un nouveau candidat :
\begin{enumerate}
    \item Vous recevez un \textbf{email de notification}
    \item Connectez-vous à votre espace
    \item Dans la section \textbf{Filleuls en attente}, cliquez sur \textbf{Valider}
    \item Le candidat pourra alors poursuivre son processus d'adhésion
\end{enumerate}

% ═══════════════════════════════════════════════════════════
\section{Module Communication}
% ═══════════════════════════════════════════════════════════

Le module Communication est une \textbf{messagerie interne} de type WhatsApp, intégrée directement dans la plateforme.

\subsection{Accéder à la Communication}

Cliquez sur \textbf{Communication} (\faBullhorn) dans la sidebar. L'interface s'affiche en \textbf{plein écran} avec :
\begin{itemize}
    \item \textbf{Panneau gauche} : liste des contacts et groupes
    \item \textbf{Panneau droit} : zone de conversation
\end{itemize}

\subsection{Types de conversations}

\begin{tabularx}{\textwidth}{llX}
\toprule
\textbf{Icône} & \textbf{Type} & \textbf{Description} \\
\midrule
\faGlobeAmericas & National & Canal de diffusion pour tous les membres \\
\faComments & Général & Discussion ouverte à tous \\
\faThLarge & Cellule & Messages aux membres de votre cellule \\
\faLayerGroup & Type de membre & Ciblage par catégorie \\
\faUsersCog & Groupe & Groupes personnalisés \\
\faUser & Individuel & Messages privés entre deux membres \\
\bottomrule
\end{tabularx}

\subsection{Envoyer un message}

\begin{etapebox}[Procédure]
\begin{enumerate}[leftmargin=1.5em]
    \item Sélectionnez un \textbf{contact} ou un \textbf{groupe} dans le panneau gauche
    \item La conversation s'ouvre dans le panneau droit
    \item Tapez votre message dans la barre en bas
    \item Appuyez sur \textbf{Entrée} ou cliquez sur le bouton \faPaperPlane
\end{enumerate}
\end{etapebox}

\subsection{Envoyer une pièce jointe}

\begin{etapebox}[Joindre un fichier]
\begin{enumerate}[leftmargin=1.5em]
    \item Cliquez sur le bouton \faPaperclip\ (trombone) dans la barre d'envoi
    \item Sélectionnez un fichier depuis votre appareil
    \item Un \textbf{aperçu} s'affiche avec le nom et la taille du fichier
    \item Ajoutez un message texte (optionnel)
    \item Cliquez sur \faPaperPlane\ pour envoyer
\end{enumerate}
\end{etapebox}

\begin{infobox}[Formats acceptés]
\textbf{Images} : JPG, PNG, GIF, WebP\\
\textbf{Documents} : PDF, Word (.doc, .docx), Excel (.xls, .xlsx), PowerPoint (.ppt, .pptx)\\
\textbf{Autres} : TXT, ZIP, RAR\\
\textbf{Taille maximale} : 10 MB par fichier
\end{infobox}

\subsection{Messages non lus}

Les messages non lus sont signalés par un \textbf{badge rouge} sur le contact dans la liste. Le badge affiche le nombre de messages non lus. Lorsque vous ouvrez une conversation, les messages sont automatiquement marqués comme \textbf{lus}.

\subsection{Rechercher un contact}

Utilisez la \textbf{barre de recherche} en haut du panneau gauche pour filtrer les contacts par nom ou email.

\subsection{Créer un groupe personnalisé}

\begin{etapebox}[Créer un groupe]
\begin{enumerate}[leftmargin=1.5em]
    \item Cliquez sur le bouton \faUsersCog\ dans la barre d'en-tête
    \item Donnez un \textbf{nom} au groupe
    \item Ajoutez une \textbf{description} (optionnel)
    \item Choisissez une \textbf{couleur} pour l'avatar
    \item Recherchez et sélectionnez les \textbf{membres} à ajouter
    \item Cliquez sur \textbf{Créer}
\end{enumerate}
\end{etapebox}

Le groupe apparaît immédiatement dans la section \textbf{Mes Groupes} de la sidebar.

% ═══════════════════════════════════════════════════════════
\section{Vidéoconférence}
% ═══════════════════════════════════════════════════════════

La plateforme intègre un système de \textbf{vidéoconférence professionnel} directement dans le navigateur, sans installation requise, sans publicité.

\subsection{Appel vidéo direct}

\begin{etapebox}[Lancer un appel vidéo]
\begin{enumerate}[leftmargin=1.5em]
    \item Ouvrez une conversation avec un contact ou un groupe
    \item Cliquez sur le bouton \faVideo\ dans la barre d'en-tête de la conversation
    \item La vidéoconférence se lance en \textbf{plein écran}
    \item Un lien de participation est \textbf{automatiquement partagé} dans le chat
    \item Les autres participants cliquent sur le lien pour rejoindre
\end{enumerate}
\end{etapebox}

\begin{astucebox}[Fonctionnalités vidéo]
\begin{itemize}[leftmargin=1.5em]
    \item \faVideo\ Caméra et \faMicrophone\ microphone activés par défaut
    \item \faDesktop\ Partage d'écran disponible
    \item \faComments\ Chat intégré dans l'appel
    \item \faHandPaper\ Lever la main
    \item \faUsers\ Nombre illimité de participants
    \item Appuyez sur \textbf{Échap} ou cliquez sur \textbf{Raccrocher} pour quitter
\end{itemize}
\end{astucebox}

\subsection{Planifier une réunion}

\begin{etapebox}[Planifier une réunion]
\begin{enumerate}[leftmargin=1.5em]
    \item Cliquez sur le bouton \faCalendarPlus\ dans la barre d'en-tête de la messagerie
    \item Remplissez les informations :
    \begin{itemize}
        \item \textbf{Titre} de la réunion (obligatoire)
        \item \textbf{Description} (optionnel)
        \item \textbf{Date et heure}
        \item \textbf{Durée} (30 min, 1h, 1h30, 2h, 3h)
    \end{itemize}
    \item Recherchez et invitez les \textbf{participants}
    \item Cliquez sur \textbf{Planifier}
    \item Un \textbf{lien de partage} est généré
\end{enumerate}
\end{etapebox}

\subsubsection{Après la planification}

\begin{itemize}
    \item \textbf{Email automatique} : tous les invités reçoivent un email avec le lien de la réunion
    \item \textbf{Lien de partage} : copiez-le et partagez-le par d'autres moyens (WhatsApp, SMS, etc.)
    \item \textbf{Rejoindre} : cliquez directement sur « Rejoindre maintenant » depuis le popup
\end{itemize}

\subsection{Consulter les réunions}

\begin{etapebox}[Voir la liste des réunions]
\begin{enumerate}[leftmargin=1.5em]
    \item Cliquez sur le bouton \faVideo\ dans la barre d'en-tête de la messagerie
    \item La liste de vos réunions s'affiche, séparées en deux sections :
    \begin{itemize}
        \item \textbf{À venir / En cours} : réunions actives avec bouton « Rejoindre »
        \item \textbf{Terminées} : historique des réunions passées
    \end{itemize}
    \item Vous pouvez \textbf{copier le lien}, \textbf{rejoindre} ou \textbf{annuler} une réunion
    \item Cliquez sur « ← Retour » pour revenir aux contacts
\end{enumerate}
\end{etapebox}

\subsection{Rejoindre via un lien}

Si quelqu'un vous envoie un lien de réunion (format \texttt{https://sebc-dushigikirane.pro/communication/?join=...}) :

\begin{enumerate}
    \item Cliquez sur le lien
    \item Connectez-vous si ce n'est pas déjà fait
    \item La vidéoconférence se lance automatiquement
\end{enumerate}

\begin{attentionbox}[Prérequis]
Pour utiliser la vidéoconférence, votre navigateur doit autoriser l'accès à votre \textbf{caméra} et votre \textbf{microphone}. Acceptez les autorisations lorsque votre navigateur les demande.
\end{attentionbox}

% ═══════════════════════════════════════════════════════════
\section{Administration (Gestionnaires)}
% ═══════════════════════════════════════════════════════════

\begin{attentionbox}[Accès restreint]
Cette section est réservée aux membres ayant le rôle \textbf{Administrateur}. Si vous êtes un membre simple, vous n'y avez pas accès.
\end{attentionbox}

\subsection{Onglets d'administration}

\begin{tabularx}{\textwidth}{lX}
\toprule
\textbf{Onglet} & \textbf{Fonctionnalités} \\
\midrule
Paramètres & Nom de l'association, frais d'adhésion, etc. \\
Pays & Gestion des pays (ajout, modification, activation) \\
Provinces & Provinces rattachées aux pays \\
Cellules & Cellules géographiques avec responsable \\
Types Ayants Droits & Types de liens familiaux \\
Types Soutien & Catégories de soutien \\
Rôles & Attribution des rôles aux membres \\
Modules & Activation/désactivation des modules et matrice d'accès \\
\bottomrule
\end{tabularx}

\subsection{Gestion des Modules}

L'onglet \textbf{Modules} permet de :
\begin{itemize}
    \item \textbf{Activer/désactiver} un module globalement
    \item \textbf{Masquer/afficher} un module dans la sidebar
    \item \textbf{Modifier} l'ordre d'affichage, l'icône, la couleur
    \item \textbf{Configurer les accès} par rôle (lire, écrire, supprimer)
\end{itemize}

% ═══════════════════════════════════════════════════════════
\section{Questions Fréquentes (FAQ)}
% ═══════════════════════════════════════════════════════════

\begin{description}[style=nextline, leftmargin=0.5cm]

\item[\textbf{Q : Je ne reçois pas le code OTP par email.}]
Vérifiez votre dossier Spam/Courrier indésirable. L'expéditeur est \texttt{link2ictgroup@gmail.com}. Si le problème persiste, attendez 2–3 minutes et redemandez un code.

\item[\textbf{Q : Mon parrain n'a pas encore validé ma candidature.}]
Depuis votre espace membre, utilisez le bouton \textbf{Relancer le parrain} pour lui envoyer un rappel par email.

\item[\textbf{Q : Je ne vois pas le module Communication dans la sidebar.}]
Assurez-vous que vous êtes connecté avec un compte approuvé. Les modules sont attribués selon votre rôle par l'administrateur.

\item[\textbf{Q : La vidéoconférence ne se lance pas.}]
Vérifiez que votre navigateur autorise l'accès à la caméra et au microphone. Essayez de rafraîchir la page. Utilisez de préférence Google Chrome.

\item[\textbf{Q : Comment supprimer un groupe que j'ai créé ?}]
Ouvrez la conversation du groupe, cliquez sur \faUserPlus\ en haut, puis sur le bouton rouge « Supprimer ».

\item[\textbf{Q : La taille de mon fichier est trop grande.}]
La limite est de \textbf{10 MB} par fichier. Compressez votre fichier ou utilisez un service de partage (Google Drive, etc.) et envoyez le lien dans le message.

\item[\textbf{Q : Comment changer mon mot de passe ?}]
Cette fonctionnalité sera implémentée prochainement. Contactez un administrateur en attendant.

\end{description}

% ═══════════════════════════════════════════════════════════
\section{Raccourcis et Astuces}
% ═══════════════════════════════════════════════════════════

\begin{tabularx}{\textwidth}{lX}
\toprule
\textbf{Action} & \textbf{Comment} \\
\midrule
Envoyer un message & Tapez le texte et appuyez sur \textbf{Entrée} \\
Joindre un fichier & Cliquez sur \faPaperclip\ puis sélectionnez le fichier \\
Lancer un appel vidéo & Cliquez sur \faVideo\ dans l'en-tête de la conversation \\
Planifier une réunion & Cliquez sur \faCalendarPlus\ dans la sidebar \\
Voir les réunions & Cliquez sur \faVideo\ dans la sidebar \\
Créer un groupe & Cliquez sur \faUsersCog\ dans la sidebar \\
Quitter un appel vidéo & Appuyez sur \textbf{Échap} ou cliquez sur « Raccrocher » \\
Rechercher un contact & Tapez dans la barre de recherche en haut \\
Retour sidebar (mobile) & Cliquez sur \faArrowLeft\ en haut de la conversation \\
\bottomrule
\end{tabularx}

% ═══════════════════════════════════════════════════════════
\section{Support}
% ═══════════════════════════════════════════════════════════

En cas de problème technique, contactez l'équipe de support :

\begin{center}
\begin{tabular}{ll}
\faEnvelope & \texttt{evanet08@gmail.com} \\
\faGlobe & \url{https://sebc-dushigikirane.pro} \\
\end{tabular}
\end{center}

\vfill
\begin{center}
\rule{0.5\textwidth}{0.5pt}\\[0.5em]
{\small\color{sebcgray} SEBC Dushigikirane — Manuel d'Utilisation v1.0 — \today}
\end{center}

\end{document}
